\section*{RESUMEN}
\addcontentsline{toc}{section}{RESUMEN}

La selección de formatos de persistencia suele realizarse por costumbre sin dimensionar la penalización computacional que impone la materialización completa en memoria mediante el modelo basado en árbol (DOM). Con el propósito de establecer criterios cuantitativos de elección para entornos con memoria acotada, se evaluó el impacto de la estructura sintáctica de CSV, JSON y XML sobre el tiempo de procesamiento y el consumo de memoria RAM en la Máquina Virtual de Java (JVM). Mediante un diseño experimental puro de corte transversal ejecutado en modo monohilo sobre OpenJDK 25, se instrumentaron $540$ corridas ($30$ repeticiones por escenario tras calentamiento del compilador JIT y filtrado por rango intercuartílico) combinando los tres formatos, dos niveles de profundidad (plano y anidado) y tres volúmenes de carga ($1$, $10$ y $50\text{ MB}$ equivalentes a hasta $325\,000$ registros). Los resultados confirmaron diferencias estadísticamente significativas ($p < 10^{-14}$, Kruskal-Wallis y ANOVA) con tamaños del efecto grandes ($\epsilon^{2} \ge 0.886$), estableciendo el orden de eficiencia $\text{CSV} > \text{JSON} > \text{XML}$. En el escenario crítico de $50\text{ MB}$ anidados, XML demandó $1\,684.39\text{ ms}$ y retuvo $629.64\text{ MB}$ en el Heap ($6.81$ veces más lento y $3.77$ veces más memoria que CSV), mientras que el anidamiento elevó el factor de expansión espacial de JSON de $5.33\times$ a $6.59\times$, igualando proporcionalmente a XML ($6.56\times$). Estas mediciones demuestran que la carga completa en árbol debe reservarse para cargas jerárquicas menores a $10\text{ MB}$ o arquitecturas con al menos $8.5$ veces el tamaño del archivo disponible en Heap.

\vspace{0.5cm}
\noindent\textbf{Palabras clave:} consumo de memoria RAM, CSV, DOM, JSON, Máquina Virtual de Java, tiempo de procesamiento, XML.